\clearpage
\section{The empirical gap certificate and midpoint-row gap checker}
\label{app:gap}
\noindent

The repository also includes a finite gap checker:
\[
  \Lean{GeneratedGapCertificate.external\_exceptionFree\_to\_certificateThreshold}.
\]
It certifies that every prime
\[
  127<r\le 191264
\]
has a midpoint-row twin witness.  So any first exception after \(127\)
must occur after \(191264\).  The Lean theorem is
\begin{itemize}[leftmargin=1.5em,itemsep=0.15em]
  \item \Lean{GeneratedGapCertificate.any\_firstException\_after\_127\_occurs\_after\_certificateThreshold}
\end{itemize}

This theorem is not a dependency of
\begin{itemize}[leftmargin=1.5em,itemsep=0.15em]
  \item \Lean{TwinPrimeExternal.arbitrarily\_large\_twins}
\end{itemize}
It is included because it records the empirical starting point of the search:
the generated finite contradiction certificate is applied beyond a verified
no-exception window after the last observed exception \(127\).

\subsection{Checker correctness}

This appendix section proves the correctness of
\[
  \Code{tools/check\_cone\_survivor\_gap.cpp}.
\]

\subsection{Specification}

The checker takes two integers \(L<V\), defaulting to
\[
  L=127,\qquad V=191264.
\]
It must verify:
\[
  \forall r,\ \Prime(r)\wedge L<r\le V
  \Rightarrow \neg\operatorname{MidpointExceptionalPrime}(r).
\]
Equivalently, for every prime \(r\in(L,V]\), it must find an index
\[
  1\le h<r
\]
such that both \(rh-1\) and \(rh+1\) are prime.

\subsection{Prime enumeration}

The function \Code{primes\_upto(V)} is the ordinary Eratosthenes sieve:
it iterates \(p=2,\ldots,V\), emits \(p\) if not previously marked composite,
and marks multiples \(p^2,p^2+p,\ldots\).  The standard invariant is:
after processing all primes \(<p\), a number \(n\ge p\) is marked if and only
if it has a prime factor \(<p\).  So at the time \(p\) is inspected, it is
unmarked if and only if it is prime.  So the returned vector is exactly
the list of primes up to \(V\).

\subsection{Reduction to indices divisible by six}

Let \(r>3\) be prime.  If \(h\) is odd, then \(rh\) is odd, so both
\[
  rh-1,\qquad rh+1
\]
are even and larger than \(2\).  So they cannot both be prime.

If \(3\nmid h\), then \(rh\not\equiv 0\pmod 3\), since \(r\ne 3\).  So one
of \(rh-1,rh+1\) is divisible by \(3\).  In the checked range \(r>127\), this
divisible neighbor is larger than \(3\), so it is composite.  So any
successful index \(h\) must be divisible by \(6\).  The checker is complete
when it searches
\[
  h=6,12,18,\ldots < r.
\]

\subsection{The twin-witness predicate}

The function \Code{cone\_survives(r,h,primes)} computes
\[
  n_- = rh-1,\qquad n_+=rh+1,
\]
and checks every prime \(q<r\).  It returns true exactly when no prime
\(q<r\) divides \(n_-\) or \(n_+\).

\begin{lemma}
Let \(r\) be prime, \(1\le h<r\), and \(n\in\{rh-1,rh+1\}\).  If no prime
\(q<r\) divides \(n\), then \(n\) is prime.
\end{lemma}

\begin{proof}
We have \(n>1\).  Also
\[
  rh+1\le r(r-1)+1=r^2-r+1<r^2,
\]
so \(n<r^2\).  If \(n\) were composite, it would have a prime divisor
\[
  q\le \sqrt n<r.
\]
This contradicts the assumption that no prime \(q<r\) divides \(n\).  So
\(n\) is prime.
\end{proof}

So, if \Code{cone\_survives(r,h,primes)} returns true, the listed \(h\)
witnesses \(\operatorname{MidpointTwinWitnessAt}(r,h)\), and \(r\) is not a
\(\operatorname{MidpointExceptionalPrime}\).

\subsection{Output correctness}

The checker loops over exactly the primes \(r\in(L,V]\).  For each such \(r\),
it searches every potentially successful index \(h\equiv 0\pmod 6\) with
\(h<r\).  If it finds a twin witness, it records a row \((r,h)\).  If not, it
records \(r\) in \Code{failedPrimes}.

The generated run produced
\[
  \Code{checkedPrimes}=17236,\qquad
  \Code{survivorRows}=17236,\qquad
  \Code{failedPrimes}=[].
\]
So every prime in \((127,191264]\) has a midpoint-row twin witness.
This proves the external Lean declaration
\[
  \Lean{GeneratedGapCertificate.external\_exceptionFree\_to\_certificateThreshold}.
\]

\subsection{Formal objects used in Appendix A}

\begin{filelist}
  \FileEntry{TwinPrimeExternal/Core.lean}
  \begin{itemize}[leftmargin=1.5em,itemsep=0.15em]
    \item \Lean{ExceptionFreeUpTo}: the finite no-exception-window predicate.
    \item \Lean{FirstExceptionAfterLastObserved}: the first-exception record used by this auxiliary checker.
    \item \Lean{firstExceptionAfterLast\_occurs\_after\_of\_exceptionFreeUpTo}: the general theorem pushing a first exception past a checked window.
    \item \Lean{firstExceptionAfterLast\_occurs\_after\_certificateThreshold\_of\_exceptionFreeUpTo}: the certificate-threshold specialization.
  \end{itemize}

  \FileEntry{TwinPrimeExternal/GeneratedGapCertificate.lean}
  \begin{itemize}[leftmargin=1.5em,itemsep=0.15em]
    \item \Lean{gapLast}: the generated last-exception bound.
    \item \Lean{gapCheckedTo}: the generated checked-to threshold.
    \item \Lean{gapSize}: the generated gap size.
    \item \Lean{checkedPrimes}: the number of checked primes in the gap run.
    \item \Lean{failedPrimes}: the generated failed-prime count.
    \item \Lean{gapLast\_eq\_lastObservedException}: checks the generated last bound.
    \item \Lean{gapCheckedTo\_eq\_certificateVerifiedTo}: checks the generated checked-to threshold.
    \item \Lean{failedPrimes\_eq\_zero}: checks that no failures were generated.
    \item \Lean{external\_exceptionFree\_to\_certificateThreshold}: the auxiliary external gap declaration.
    \item \Lean{any\_firstException\_after\_127\_occurs\_after\_certificateThreshold}: derives the advertised first-exception lower bound.
  \end{itemize}
\end{filelist}
